\documentclass[11pt]{article}
\usepackage[margin=1in]{geometry}
\usepackage{amsmath, amssymb, amsthm, bm, mathtools}
\usepackage{booktabs, array}
\usepackage{enumitem}
\usepackage{siunitx}
\usepackage{hyperref}
\hypersetup{colorlinks=true, linkcolor=blue, urlcolor=blue, citecolor=blue}
\sisetup{round-mode=places,round-precision=3,group-separator={,}}

\newcommand{\dd}{\,\mathrm{d}}
\newcommand{\R}{\mathbb{R}}


\title{Solutions -- Practice Problem Set 1}
\author{ECON 308 -- International Economic Relations}
\date{Fall 2025}

\begin{document}
\maketitle



\begin{enumerate}[leftmargin=*]
\item \textbf{Gravity intuition.} Larger GDPs imply more production and income, scaling the number of potential buyers and sellers; microfoundations include love-of-variety models and market-size effects that raise the mass of matches. Distance proxies trade costs (freight, time-to-market, information frictions), which depress the intensive and extensive margins of trade. Search/information models also imply distance lowers match probabilities, while per-unit transport costs directly reduce delivered profitability.

\item \textbf{Home bias and border effects.} Home bias is the empirical tendency for domestic goods to be consumed disproportionately relative to foreign goods even after controlling for observables. A border effect can be detected by including a same-country dummy in a gravity regression; a large, positive coefficient implies much higher domestic trade than predicted by distance/size alone. Interpretation: additional unobserved frictions (legal regimes, standards, currencies) or preferences (brand/home familiarity).

\item \textbf{Openness vs.\ size.} Small economies typically have narrower domestic markets and fewer varieties, pushing firms to export to exploit scale and consumers to import for variety. Large economies internalize more trade within borders (many regions/sectors), mechanically lowering $(X{+}M)/\text{GDP}$. Fragmented supply chains also concentrate cross-border flows in small open economies.

\item \textbf{Inter- vs.\ intra-industry trade.} Inter-industry trade exchanges different goods across industries (e.g., crude oil for machinery). Intra-industry trade exchanges similar varieties within the same industry (e.g., Germany and France trading differentiated automobiles), often driven by product differentiation and economies of scale.
\end{enumerate}




\begin{enumerate}[leftmargin=*,resume]
\item \textbf{Absolute vs.\ comparative advantage.} Absolute advantage means using fewer inputs per unit of output; comparative advantage means having a lower \emph{opportunity cost}. A country with absolute disadvantage in both goods can still gain from trade by specializing where its disadvantage is smallest (lower OC). Relative prices then reflect technology differences, allowing both countries to consume beyond autarky PPFs.

\item \textbf{Relative wages under free trade.} With perfect competition and specialization, each country's wage equals the value of marginal product in the sector it produces: $w=\dfrac{P_g}{a_{Lg}}$. The \emph{relative} wage across countries is pinned down by world prices and comparative advantage: the high-wage country must have sufficiently high productivity in its export good to remain competitive at world prices.
\newpage
\item \textbf{PPF vs.\ CPF.} The PPF is technologically feasible output combinations; the CPF reflects what a country can \emph{consume} given world prices and its export revenue. Under autarky CPF $=$ PPF; with trade, selling at world prices and buying imports typically lets the CPF lie strictly outside the PPF, delivering gains from trade.

\item \textbf{Unit labor costs and prices.} In autarky, competitive pricing implies $P_g=w\,a_{Lg}$, so relative prices equal unit labor cost ratios. With trade, world prices need not equal a country's own unit costs; instead, specialization ensures the country produces the good where it has the lowest unit cost at world prices, and wages adjust to equalize unit costs in the produced sector.
\end{enumerate}




\begin{enumerate}[leftmargin=*, resume]
\item \textbf{Gravity Model: Baseline Predictions and Elasticities.}

\emph{Given:} $T_{ij}=10^{-6}\dfrac{Y_iY_j}{D_{ij}}$, with $(Y_A,Y_B,Y_C)=(2000,1000,500)$ (billions), and distances in km:
$D_{AB}=6{,}000$, $D_{AC}=1{,}500$, $D_{BC}=2{,}000$.

\begin{enumerate}
\item \emph{Predicted bilateral exports.}
\[
\begin{aligned}
T_{AB}&=10^{-6}\frac{(2000)(1000)}{6000}
=10^{-6}\cdot \frac{2{,}000{,}000}{6000}
=10^{-6}\cdot 333.\overline{3}
=0.000333\ \text{(billions)}\\
&\Rightarrow \ \ \$0.333\ \text{million.}\\[6pt]
T_{AC}&=10^{-6}\frac{(2000)(500)}{1500}
=10^{-6}\cdot \frac{1{,}000{,}000}{1500}
=10^{-6}\cdot 666.\overline{6}
=0.000667\ \text{(billions)}\\
&\Rightarrow \ \ \$0.667\ \text{million.}\\[6pt]
T_{BC}&=10^{-6}\frac{(1000)(500)}{2000}
=10^{-6}\cdot \frac{500{,}000}{2000}
=10^{-6}\cdot 250
=0.000250\ \text{(billions)}\\
&\Rightarrow \ \ \$0.250\ \text{million.}
\end{aligned}
\]

\item \emph{Effect of a 10\% fall in distance on $T_{AB}$.} The log-specification gives the elasticity of trade with respect to distance as $-\gamma=-1$. Exactly:
\[
\frac{T_{AB}^{\,\text{new}}}{T_{AB}^{\,\text{old}}}
=\left(\frac{D_{AB}^{\text{new}}}{D_{AB}^{\text{old}}}\right)^{-\gamma}
=(0.9)^{-1}=\frac{1}{0.9}=1.\overline{1}
\]
So $T_{AB}$ rises by $11.\overline{1}\%$.

\item \emph{Effect of 5\% GDP increases in both A and B.} Because $T_{ij}\propto Y_iY_j$,
\[
\frac{T_{AB}^{\,\text{new}}}{T_{AB}^{\,\text{old}}}=(1.05)(1.05)=1.1025 \Rightarrow \text{increase of } 10.25\%.
\]

\item \emph{Log residual against observed \$1.3 bn.} 
\[
\ln T^{\text{obs}}_{AB}-\ln T^{\text{pred}}_{AB}
=\ln(1.3)-\ln(0.000333\ldots)\approx 8.269.
\]
A positive residual indicates $AB$ trade exceeds gravity prediction, consistent with omitted frictions/facilitators (e.g., common language, trade agreement, high-quality transport links, or substantial services trade not captured by simple distance).
\end{enumerate}

\newpage
\item \textbf{Openness: Trade-to-GDP Ratios.}

\emph{Data (billions; population in millions).}
\[
\begin{array}{@{}lrrrr@{}}
\toprule
\text{Country} & X & M & \text{GDP} & \text{Pop.}\\
\midrule
D & 120 & 100 & 1000 & 50\\
E & 400 & 350 & 2500 & 200\\
F & 70 & 110 & 600 & 10\\
\bottomrule
\end{array}
\]

\begin{enumerate}
\item Openness $(X{+}M)/\text{GDP}$:
\[
\begin{aligned}
\text{D: }&\frac{220}{1000}=0.22=22\%,\quad
\text{E: }\frac{750}{2500}=0.30=30\%,\quad
\text{F: }\frac{180}{600}=0.30=30\%.
\end{aligned}
\]

\item Per-capita trade (domestic currency per person).\footnote{(billions)/(millions) $=\$1000$ per person.}
\[
\begin{aligned}
\text{D: }&X/Pop = \frac{120}{50}=2.4 \Rightarrow 2{,}400,\quad
M/Pop = \frac{100}{50}=2.0 \Rightarrow 2{,}000;\\
\text{E: }&X/Pop = \frac{400}{200}=2.0 \Rightarrow 2{,}000,\quad
M/Pop=\frac{350}{200}=1.75 \Rightarrow 1{,}750;\\
\text{F: }&X/Pop = \frac{70}{10}=7.0 \Rightarrow 7{,}000,\quad
M/Pop=\frac{110}{10}=11.0 \Rightarrow 11{,}000.
\end{aligned}
\]

\item Ranking by openness: E = F (30\%) $>$ D (22\%). Small economies often import/export a larger share of GDP due to limited domestic variety and scale---they rely on trade to access varieties and exploit economies of scale.
\end{enumerate}

\item \textbf{Borders and Common Language Dummies in Gravity.}

Adjustment: multiply baseline $T_{ij}$ by $(1+\lambda_{CL})$ if common language (CL) and by $(1-\lambda_{BOR})$ if a border friction applies. With $\lambda_{CL}=0.20$, $\lambda_{BOR}=0.30$ and all pairs separated by a border, while only $AB$ shares CL:
\[
\begin{aligned}
T_{AB}^{\text{adj}}&=T_{AB}\times (1{+}0.20)\times(1{-}0.30)
=0.000333\ldots \times 1.2 \times 0.7
=0.000280,\\
T_{AC}^{\text{adj}}&=T_{AC}\times 0.7
=0.000667\times 0.7
=0.000467,\\
T_{BC}^{\text{adj}}&=T_{BC}\times 0.7
=0.000250\times 0.7
=0.000175.
\end{aligned}
\]
\emph{Language effect percent increase for $AB$:} holding the border at $0.7$, the CL multiplies trade by $1.2$, i.e., $+20\%$.
\emph{Border friction that exactly offsets CL:} solve $(1+\lambda_{CL})(1-\lambda_{BOR})=1$:
\[
1.2\,(1-\lambda_{BOR})=1 \ \Rightarrow\ 1-\lambda_{BOR}=\tfrac{1}{1.2}=\tfrac{5}{6} \ \Rightarrow\ \lambda_{BOR}=\tfrac{1}{6}\approx 0.167.
\]

\item \textbf{Growth and Gravity Over Time.}

\emph{Given:} annual real growth $g_A=2.5\%$, $g_C=4\%$, 10 years.
\[
Y_A(10)=2000(1.025)^{10}\approx 2560.169,\quad
Y_C(10)=500(1.04)^{10}\approx 740.122.
\]
With $D_{AC}$ unchanged and $\beta_1=\beta_2=\gamma=1$,
\[
\frac{T_{AC}(10)}{T_{AC}(0)}
=\frac{Y_A(10)}{Y_A(0)}\cdot \frac{Y_C(10)}{Y_C(0)}
=(1.025)^{10}(1.04)^{10}\approx 1.895,
\]
i.e., an $89.5\%$ increase is predicted. If observed growth is only $12\%$, shortfalls could reflect higher trade costs (tariffs/NTBs), weaker logistics, currency/financial frictions, or policy shocks (sanctions, standards) that the simple gravity baseline omits.
\end{enumerate}





\newpage

\begin{enumerate}[leftmargin=*, resume]
\item \textbf{Two-Country, Two-Good Ricardian Basics.}

\emph{Given:} Unit labor requirements
\[
\begin{array}{@{}lcc@{}}
\toprule
 & a_{LW} & a_{LC}\\
\midrule
H & 2 & 4\\
F & 6 & 3\\
\bottomrule
\end{array}
\qquad L_H=L_F=240.
\]

\begin{enumerate}
\item \emph{Opportunity costs and comparative advantage.} In the Ricardian model,
\[
\text{OC}_H(W \text{ in } C)=\frac{a_{LW}}{a_{LC}}=\frac{2}{4}=0.5,\quad
\text{OC}_H(C \text{ in } W)=\frac{a_{LC}}{a_{LW}}=\frac{4}{2}=2.
\]
\[
\text{OC}_F(W \text{ in } C)=\frac{6}{3}=2,\quad
\text{OC}_F(C \text{ in } W)=\frac{3}{6}=0.5.
\]
Home (H) has lower OC in $W$ ($0.5<2$), so $H$ has comparative advantage in $W$; Foreign (F) in $C$.

\item \emph{PPFs and autarky prices.} PPFs satisfy $a_{LW}W+a_{LC}C\le L$.
\[
\text{H: } 2W+4C\le 240 \Rightarrow W=120-2C\ (\text{slope }-2).
\qquad
\text{F: } 6W+3C\le 240 \Rightarrow W=40-\tfrac{1}{2}C\ (\text{slope }-\tfrac{1}{2}).
\]
Under perfect competition, autarky relative price equals the unit-cost ratio:
\[
\left(\frac{P_W}{P_C}\right)^{A}_H=\frac{a_{LW}}{a_{LC}}=\tfrac{1}{2}=0.5,
\qquad
\left(\frac{P_W}{P_C}\right)^{A}_F=\frac{6}{3}=2.
\]

\item \emph{Range for complete specialization.} With $H$ in $W$ and $F$ in $C$,
\[
0.5 < \frac{P_W}{P_C} < 2.
\]

\item \emph{Specialization at $P_W/P_C=1.2$.} Since $0.5<1.2<2$, $H$ specializes in $W$, $F$ in $C$.
\[
\text{H output: } W_H=\frac{L_H}{a_{LW}}=\frac{240}{2}=120,\quad C_H=0.
\qquad
\text{F output: } C_F=\frac{L_F}{a_{LC}}=\frac{240}{3}=80,\quad W_F=0.
\]
$H$ exports $W$ and imports $C$; $F$ the reverse.

\item \emph{Real wages under free trade.} With specialization, $w_H=\dfrac{P_W}{a_{LW}}$ and $w_F=\dfrac{P_C}{a_{LC}}$.
\[
\text{H:}\quad \frac{w_H}{P_W}=\frac{1}{a_{LW}}=\frac{1}{2}=0.5\ \text{(units of $W$)},\quad
\frac{w_H}{P_C}=\frac{P_W/P_C}{a_{LW}}=\frac{1.2}{2}=0.6\ \text{(units of $C$)}.
\]
\[
\text{F:}\quad \frac{w_F}{P_C}=\frac{1}{a_{LC}}=\frac{1}{3}\ \text{(units of $C$)},\quad
\frac{w_F}{P_W}=\frac{1}{a_{LC}}\cdot\frac{P_C}{P_W}=\frac{1}{3\cdot 1.2}=\frac{1}{3.6}\approx 0.278\ \text{(units of $W$)}.
\]
\end{enumerate}

\newpage

\item \textbf{Terms of Trade and Real Incomes (price falls to $P_W/P_C=0.8$).}
\begin{enumerate}
\item Since $0.5<0.8<2$, specialization pattern is unchanged: $H\to W$, $F\to C$.

\item \emph{Real wages at $P_W/P_C=0.8$.}
\[
\text{H:}\ \frac{w_H}{P_W}=\frac{1}{2}=0.5\ \text{W},\quad
\frac{w_H}{P_C}=\frac{0.8}{2}=0.4\ \text{C}.
\quad
\text{F:}\ \frac{w_F}{P_C}=\frac{1}{3}\ \text{C},\quad
\frac{w_F}{P_W}=\frac{1}{3\cdot 0.8}=\frac{1}{2.4}\approx 0.417\ \text{W}.
\]
\emph{Comparison to autarky:} Autarky real wages are $H:(1/2\ \text{W},\ 1/4\ \text{C})$; $F:(1/6\ \text{W},\ 1/3\ \text{C})$. Both countries gain from trade relative to autarky.

\item \emph{Terms of trade (TOT) effects.} $H$ exports $W$, so its TOT is $P_W/P_C$. When this falls from $1.2$ to $0.8$, $H$'s purchasing power in $C$ falls ($0.6\to 0.4$), while $F$'s TOT $=(P_C/P_W)$ improves (its real wage in $W$ rises $0.278\to 0.417$).
\end{enumerate}

\item \textbf{World Relative Supply with Three Countries.}

\emph{Given:}
\[
\begin{array}{@{}lcc@{}}
\toprule
 & a_{LW} & a_{LC}\\
\midrule
1 & 1 & 2\\
2 & 2 & 3\\
3 & 4 & 5\\
\bottomrule
\end{array}
\qquad L_1=L_2=L_3=100.
\]
Compute $a_{LW}/a_{LC}$: Country 1: $0.5$; 2: $2/3\approx 0.667$; 3: $4/5=0.8$.

\begin{enumerate}
\item \emph{Ordering by comparative advantage in $W$:} $1$ (best), $2$, $3$.

\item \emph{World relative supply (WRS).} Each country produces $W$ if $P_W/P_C\ge a_{LW}/a_{LC}$; otherwise $C$.
\[
\begin{array}{l|c|c|c}
\text{Price range} & \text{W producers} & \text{W supply} & \text{C supply}\\
\hline
P<0.5 & \text{none} & 0 & \frac{100}{2}+\frac{100}{3}+\frac{100}{5}=50+33.\overline{3}+20=103.\overline{3}\\
0.5<P<2/3 & 1 & \frac{100}{1}=100 & \frac{100}{3}+\frac{100}{5}=33.\overline{3}+20=53.\overline{3}\\
2/3<P<0.8 & 1,2 & \frac{100}{1}+\frac{100}{2}=150 & \frac{100}{5}=20\\
P>0.8 & 1,2,3 & \frac{100}{1}+\frac{100}{2}+\frac{100}{4}=175 & 0
\end{array}
\]
Thus the step values of $RS^{W/C}\equiv W/C$ are: $0$; $\dfrac{100}{53.\overline{3}}\approx 1.875$; $\dfrac{150}{20}=7.5$; and $+\infty$.
Kinks at $P=0.5,\ 2/3,\ 0.8$.

\item \emph{Equilibrium with $RD^{W/C}=1.1-0.2(P_W/P_C)$.}
Because $RD$ is continuous and decreasing, and $RS$ is a step function, the market clears either (i) on a step (if $RD$ hits a step value), or (ii) at a kink with the marginal country partially specialized. Evaluate $RD$ at kinks:
\[
RD(0.5)=1.1-0.2(0.5)=1.0,\quad RD(2/3)\approx 0.967,\quad RD(0.8)=0.94.
\]
All are \emph{below} the adjacent step value $1.875$. Hence equilibrium occurs at the \emph{lowest} kink $P_W/P_C=0.5$, with Country 1 splitting labor between $W$ and $C$ to make $RS=RD=1$.

Let $x\in[0,1]$ be the share of Country 1's labor in $W$ at $P=0.5$.
\[
W=\frac{100x}{1}=100x,
\qquad
C=\frac{100(1-x)}{2}+\frac{100}{3}+\frac{100}{5}=50(1-x)+53.\overline{3}=103.\overline{3}-50x.
\]
Set $W/C=1$:
\[
100x=103.\overline{3}-50x \ \Rightarrow\ 150x=103.\overline{3}\ \Rightarrow\ x\approx 0.6889.
\]
\emph{Production at equilibrium:} Country 1 uses $\approx 68.9\%$ of $L_1$ in $W$; Countries 2 and 3 produce only $C$.
\end{enumerate}

\item \textbf{Productivity Shock in $C$ at Home ($a_{LC}:4\to 3$).}

\begin{enumerate}
\item \emph{New opportunity cost at Home.}
\[
\text{OC}_H(W \text{ in } C)=\frac{a_{LW}}{a_{LC}^{\,\text{new}}}=\frac{2}{3}\approx 0.667
\quad\big(\text{was }0.5\big).
\]
Home becomes relatively \emph{better} at $C$ (lower unit cost), so the complete-specialization range with $H\to W, F\to C$ tightens to
\[
\frac{P_W}{P_C}\in\left(\frac{2}{3},\,2\right).
\]

\item \emph{Effect on world relative supply and price.} Improving $H$'s productivity in $C$ shifts world relative supply of $W$ \emph{leftward} (for a given price, relatively less $W$ supplied), pushing the equilibrium $\dfrac{P_W}{P_C}$ \emph{upward}.

\item \emph{Real wages vs.\ pre-shock free trade.} With $H$ still specializing in $W$, $w_H=\dfrac{P_W}{2}$ (unchanged technology in $W$). Thus
\[
\frac{w_H}{P_W}=\frac{1}{2}\ \text{(unchanged in $W$)},\qquad
\frac{w_H}{P_C}=\frac{1}{2}\cdot\frac{P_W}{P_C}\ \text{(rises as $P_W/P_C$ rises)}.
\]
Foreign continues in $C$ with $w_F=\dfrac{P_C}{3}$, so
\[
\frac{w_F}{P_C}=\frac{1}{3}\ \text{(unchanged)},\qquad
\frac{w_F}{P_W}=\frac{1}{3}\cdot\frac{P_C}{P_W}=\frac{1}{3}\cdot\frac{1}{P_W/P_C}\ \text{(falls)}.
\]
Hence Home's terms of trade improve and its real wage in $C$ rises; Foreign loses in terms of $W$. Overall, the country experiencing productivity growth (Home) captures a larger share of the gains.
\end{enumerate}
\end{enumerate}




\end{document}
